\documentclass[11pt]{article}
\usepackage[a4paper,margin=26mm]{geometry}
\usepackage{amsmath,amssymb}
\setlength{\parindent}{0pt}
\setlength{\parskip}{7pt}
\emergencystretch=2em
\title{Pairwise negative correlation does not ensure feasibility\\
\large Disproof of conjecture 00000008790}
\author{ziangni-sys}
\date{4 October 2026}
\begin{document}
\maketitle
The last clause of the English conjecture asserts that pairwise negatively
correlated rounding is always feasible over any down-closed constraint
family. The Chinese last clause states the same universal assertion.
We disprove that clause by a marginal-preserving random rounding with
strictly negative correlations. This concerns sufficiency of the correlation
property; we do not claim a failure of a particular pipage implementation.

\textbf{The constraint and fractional solution.}
Let the ground set be $\{1,2\}$ and let
\[
 \mathcal F=\{\varnothing,\{1\},\{2\}\}
          =\{S\subseteq\{1,2\}:|S|\leq1\}.
\]
This family is down-closed: a subset cannot have larger cardinality.
Its integer feasible points are zero-one vectors with $y_1+y_2\leq1$.
The LP relaxation
\[
 0\leq x_1,x_2\leq1,\qquad x_1+x_2\leq1
\]
has the feasible point $x=(1/2,1/2)$. This is also the average of the
feasible integer vectors $(1,0)$ and $(0,1)$ and belongs to their convex hull.

\textbf{The actual rounding distribution.}
On $\Omega=\{0,1\}^2$, assign these probability masses:
\[
\begin{array}{c|rrrr}
\omega &(0,0)&(0,1)&(1,0)&(1,1)\\ \hline
\mathbb P(\{\omega\})&1/8&3/8&3/8&1/8 .
\end{array}
\]
They are nonnegative and sum to one. Set $Y_i(\omega)=\omega_i$ and
$S(\omega)=\{i:Y_i(\omega)=1\}$. Both marginals satisfy
\[
 \mathbb E Y_i=\mathbb P(Y_i=1)=\tfrac12=x_i.
\]
For the only distinct coordinate pair, in either order,
\[
 \mathbb P(Y_1=Y_2=1)=\tfrac18
       <\tfrac14=\mathbb P(Y_1=1)\mathbb P(Y_2=1),
\]
and also
\[
 \mathbb P(Y_1=Y_2=0)=\tfrac18
       <\tfrac14=\mathbb P(Y_1=0)\mathbb P(Y_2=0).
\]
It satisfies both selected-coordinate and complement-coordinate pairwise
negative correlation. Its covariance is $1/8-1/4=-1/8$.

\newpage
\textbf{Positive probability of infeasibility.}
The output $S$ is outside $\mathcal F$ exactly at $\omega=(1,1)$.
Consequently
\[
 \mathbb P(S\notin\mathcal F)=\tfrac18>0.
\]
Feasibility does not hold almost surely. Down-closedness and pairwise
negative correlation, with exact feasible marginals and negative
correlation of complements, do not imply universal feasibility.

A procedure enforcing the constraint can avoid this output, for example
by giving probability $1/2$ each to $(1,0)$ and $(0,1)$. Our example concerns
the sufficiency of pairwise negative correlation alone. It makes no
separate assertion about concrete pipage algorithms, approximation
ratios, deterministic constructions, or perturbation radii. One false
universal conjunct disproves the stated conjecture.

\textbf{Formal correspondence.}
The Lean project constructs a finite probability mass function with
\texttt{PMF.ofFintype}, proves normalization, and uses its actual
associated probability measure. The definitions \texttt{bit},
\texttt{selected}, \texttt{family}, and \texttt{feasibleFractional}
encode the zero-one coordinates, random output set, cardinality
constraint, and displayed LP inequalities.

Probabilities are the real values of the actual measure of the events.
\texttt{zero\_one} proves that the coordinates are indicators.
\texttt{expected\_coordinate} computes their finite weighted expectations.
\texttt{marginals} and \texttt{complements} compute event probabilities
for every coordinate. \texttt{joint\_selected} and
\texttt{joint\_unselected} compute joint probabilities for every distinct
ordered pair. \texttt{pairwise\_negative} verifies both inequalities.

\texttt{down\_closed} proves the general subset property;
\texttt{fractional\_feasible} checks every LP inequality.
\texttt{infeasible\_probability} evaluates the actual event that the
output violates the cardinality bound. The final \texttt{counterexample}
assembles these hypotheses and the strictly positive failure probability.

\textbf{Reproduction.}
Use Lean 4.19.0 and public Mathlib revision
\begin{center}\small\texttt{c44e0c8ee63ca166450922a373c7409c5d26b00b}\end{center}
Run \texttt{lake build} inside \texttt{lean/}.
The source prints axiom audits for the main results. No unproved
certificates, custom axioms, or computational trust extensions are used.
The accompanying notes record the completed build and PDF validation.
\end{document}
